% TC03 — บทบาทเพิ่มเติมและการเข้าถึงนิสิต (IDOR)
\begin{testscenario}{3}{3}
  \tcmeta{TC03}{CS 2568/01}{ระบบเว็บแอปพลิเคชัน ``NUCARE''}{นายสมชาย ใจดี}{วันที่ 22 ตุลาคม พ.ศ.\ 2569}{1.0}
  \tcname{การตรวจสอบสิทธิ์บทบาทเพิ่มเติมและการเข้าถึงข้อมูลนิสิต}
  \tcpre{มีหน่วยทดสอบ \tccode{roles.test.ts} และ \tccode{access-control.test.ts}
    ในเซิร์ฟเวอร์ และรันได้ด้วย \tccode{npm run test:unit}}
  \tcdesc{ทดสอบว่าการตรวจบทบาทรองรับทั้งบทบาทหลักและบทบาทเพิ่มเติม
    และการเข้าถึงนิสิตรายคนถูกจำกัดตามการมอบหมาย (ป้องกัน IDOR) สำหรับที่ปรึกษา/ผู้สอน}

  \tccase{1}{ตรวจสอบ hasAssignedRole เมื่อมีบทบาทเพิ่มเติมใน payload}{%
    \begin{tcsteps}
      \item เปิดไฟล์ \tccode{src/lib/roles.test.ts}
      \item รันกรณี RL-02 ด้วยผู้ใช้ role=ADVISOR และ roles รวม ADMIN
      \item ตรวจว่า \tccode{hasAssignedRole} สำหรับ ADMIN เป็นจริง
    \end{tcsteps}}{%
    ระบบนับทั้งบทบาทหลักและรายการ roles แล้วคืนค่า true สำหรับ ADMIN}{%
    รัน test:unit ผ่านกรณี RL-02 ตามที่คาดหวัง}{Pass}

  \tccase{2}{ตรวจสอบ canAccessStudent ปฏิเสธที่ปรึกษาที่ไม่มี assignment}{%
    \begin{tcsteps}
      \item เปิดไฟล์ \tccode{src/lib/access-control.test.ts}
      \item ตั้ง mock ให้ findActiveAdvisorAssignment คืนค่า false
      \item เรียก canAccessStudent ของที่ปรึกษาที่ไม่มี assignment
    \end{tcsteps}}{%
    ฟังก์ชันคืนค่า false และไม่อนุญาตเข้าถึงนิสิตที่ไม่อยู่ในความดูแล}{%
    รัน test:unit ผ่านกรณี AC-03 ตามที่คาดหวัง}{Pass}

  \tccase{3}{ตรวจสอบผู้บริหารไม่เข้าถึงนิสิตรายคนผ่าน canAccessStudent}{%
    \begin{tcsteps}
      \item เรียก canAccessStudent ด้วยบทบาท EXECUTIVE ในหน่วยทดสอบ AC-05
      \item ตรวจค่าที่คืน
    \end{tcsteps}}{%
    ฟังก์ชันคืนค่า false ตามนโยบาย PDPA / ไม่เปิดประวัติรายบุคคล}{%
    รัน test:unit ผ่านกรณี AC-05 ตามที่คาดหวัง}{Pass}
\end{testscenario}
